\documentclass[11pt]{article}

\usepackage[a4paper,margin=1in]{geometry}
\usepackage{amsmath,amssymb}
\usepackage{graphicx}
\usepackage{booktabs}
\usepackage{hyperref}
\usepackage{enumitem}
\usepackage{float}
\usepackage{caption}

\graphicspath{{./figs/}}

\title{CE 488 / 5588 \\ Multi-Class Classification -- History and Trends \\ \large Case Study: Handwritten Digit Recognition}
\author{}
\date{}

\begin{document}

\maketitle
% \tableofcontents
\newpage

% =========================================================
\section{Introduction}

Handwritten digit recognition is one of the most influential benchmark problems in machine learning. 
Although visually simple, it encapsulates nearly all core elements of supervised learning:

\begin{itemize}
    \item Multi-class classification (10 classes),
    \item High-dimensional input space (pixel grids),
    \item Nonlinear decision boundaries,
    \item Representation learning,
    \item Statistical generalization and robustness challenges.
\end{itemize}

The problem serves as a controlled experimental testbed for understanding:
\begin{enumerate}
    \item Classical machine learning pipelines,
    \item Deep neural network evolution,
    \item Error bounds and Bayes limits,
    \item Overfitting and distribution shift.
\end{enumerate}

Despite being considered ``solved,'' MNIST remains pedagogically powerful for understanding how AI systems learn and where their limits lie.

% =========================================================
\section{Dataset Foundations}

\subsection{NIST Special Database 19 (SD19)}

The foundation dataset for modern handwritten character recognition is the 
\textbf{NIST Special Database 19 (SD19)}.

\begin{itemize}
    \item Collected from approximately 3{,}600 writers.
    \item 814{,}255 segmented character images.
    \item Image resolution: $128 \times 128$ grayscale.
    \item 62 ASCII classes:
    \begin{itemize}
        \item 10 digits,
        \item 26 uppercase letters,
        \item 26 lowercase letters.
    \end{itemize}
\end{itemize}

This dataset reflects significant diversity in handwriting styles and stroke variability.

\subsection{MNIST Benchmark Dataset}

The MNIST dataset is a curated subset derived from SD19, designed for controlled benchmarking.

\begin{itemize}
    \item 60{,}000 training samples,
    \item 10{,}000 test samples,
    \item $28 \times 28$ grayscale images,
    \item 10 digit classes: $\{0,1,\dots,9\}$,
    \item Size-normalized and centered.
\end{itemize}

MNIST became the canonical benchmark for multi-class classification for over two decades.

\subsection{Sample MNIST Digits}

\begin{figure}[H]
\centering
\includegraphics[width=0.85\textwidth]{MnistExamplesModified.png}
\caption{Example MNIST handwritten digits. Observe variations in stroke thickness, skew, and ambiguity between certain classes (e.g., 4 vs 9).}
\end{figure}

Even at this small scale, class ambiguity is evident.

% =========================================================
\section{Formal Multi-Class Classification Setup}

Let the test dataset be:

\[
D_{\text{test}} = \{(x_i, y_i)\}_{i=1}^N,
\]

where:
\begin{itemize}
    \item $x_i \in \mathbb{R}^{784}$ is the flattened $28\times 28$ image,
    \item $y_i \in \{0,1,\dots,9\}$ is the true label.
\end{itemize}

A classifier $f$ produces predictions:

\[
\hat{y}_i = f(x_i).
\]

\subsection{Error Rate and Accuracy}

The empirical error rate is defined as:

\[
\epsilon =
\frac{1}{N}
\sum_{i=1}^N
\mathbf{1}\{\hat{y}_i \neq y_i\}.
\]

Accuracy is:

\[
\text{Accuracy} = 1 - \epsilon.
\]

This corresponds to 0--1 loss evaluation.

% =========================================================
\section{Evolution of Benchmark Performance}

\subsection{Milestone Models}

\begin{itemize}
    \item \textbf{1998 -- LeNet-5 (Convolutional Neural Network)} \\
    Error rate: $\epsilon = 0.95\%$.

    \item \textbf{2011 -- Multi-Column Deep Neural Network (MCDNN)} \\
    Error rate: $\epsilon = 0.23\%$.

    \item \textbf{2013 -- DropConnect CNN} \\
    Error rate: $\epsilon = 0.21\%$.

    \item \textbf{2017 -- Capsule Networks} \\
    Error rate: $\epsilon = 0.16\%$.
\end{itemize}

\subsection{Common Characteristics}

All top-performing approaches share:

\begin{itemize}
    \item End-to-end training directly on pixels,
    \item Minimal preprocessing,
    \item Hierarchical feature learning,
    \item Joint optimization of representation and classifier.
\end{itemize}

Feature extraction is internal to the model.

% =========================================================
\section{Classical Feature-Based Methods}

Before deep learning dominance, pipelines were typically:

\[
\text{Image} \rightarrow \text{Feature Extraction} \rightarrow \text{Classifier}.
\]

\subsection{HOG + SVM}

\begin{itemize}
    \item Histogram of Oriented Gradients (HOG),
    \item Support Vector Machine classifier,
    \item Error rate $\approx 0.64\%$.
\end{itemize}

\subsection{Random Forests}

\begin{itemize}
    \item Often combined with HOG features,
    \item Error rate $\approx 2\%$--$3\%$.
\end{itemize}

\subsection{Naïve Bayes}

\begin{itemize}
    \item Gaussian likelihood per pixel,
    \item Independence assumption among pixels,
    \item Error rate $\approx 4\%$--$15\%$.
\end{itemize}

These methods separate representation learning from classification.

% =========================================================
\section{Can We Reach 100\% Accuracy?}

The best reported test error on MNIST is approximately:

\[
\epsilon \approx 0.16\%.
\]

This raises a fundamental question:

Is 100\% accuracy possible?

% =========================================================
\section{Bayes Error and Irreducible Uncertainty}

Define the Bayes error rate:

\[
\varepsilon_{\text{Bayes}} = \min_f \Pr(f(x) \neq y).
\]

Even with infinite data and perfect modeling, no classifier can outperform this limit.

Estimated Bayes error for MNIST:

\[
\varepsilon_{\text{Bayes}} \approx 0.1\% - 0.2\%.
\]

\subsection{Sources of Irreducible Error}

\begin{itemize}
    \item Ambiguous handwriting,
    \item Mislabeling,
    \item Overlapping class manifolds,
    \item Intrinsic noise.
\end{itemize}

Some digits are fundamentally indistinguishable.

\bigskip
\textbf{Key Insight:} Not all error is model error; some error is information-theoretic.

% =========================================================
\section{Beyond Testing: Generalization}

High test accuracy does not guarantee:

\begin{itemize}
    \item Robustness,
    \item Out-of-distribution performance,
    \item Stability under perturbation.
\end{itemize}

\subsection{Distribution Shift}

MNIST assumes:

\begin{itemize}
    \item Same scanning pipeline,
    \item Same normalization,
    \item Same centering,
    \item Same resolution.
\end{itemize}

Real-world deployment may involve:

\begin{itemize}
    \item Different lighting,
    \item Camera distortions,
    \item Writer demographics,
    \item Image corruption.
\end{itemize}

\subsection{Adversarial Vulnerability}

Small perturbations to pixel intensities can cause confident misclassification.

Thus, robustness must be evaluated separately from raw accuracy.

% =========================================================
\section{Deeper Conceptual Lessons}

Digit recognition embodies:

\begin{itemize}
    \item High-dimensional nonlinear manifolds,
    \item Bias--variance tradeoff,
    \item Statistical vs. information-theoretic limits,
    \item Distributional robustness challenges.
\end{itemize}

These principles extend directly to:

\begin{itemize}
    \item Medical imaging,
    \item Structural damage detection,
    \item Infrastructure monitoring,
    \item Autonomous systems.
\end{itemize}

The mathematical structure remains the same; only the domain changes.

% =========================================================
\section{Suggested Exercises}

\begin{enumerate}
    \item Train logistic regression on MNIST.
    \item Compare linear and kernel SVM performance.
    \item Add Gaussian noise and re-evaluate accuracy.
    \item Rotate digits by $\pm 10^\circ$ and test.
    \item Train a CNN and compare robustness.
    \item Compute and interpret the confusion matrix.
\end{enumerate}

% =========================================================
\section{Conclusion}

MNIST is more than a solved benchmark.

It provides a controlled environment to understand:

\begin{itemize}
    \item Representation learning,
    \item End-to-end optimization,
    \item Bayes limits,
    \item Generalization versus memorization,
    \item Robustness beyond benchmark performance.
\end{itemize}

The deeper question is not whether we can reach 100\% accuracy, but whether the model truly understands the concept of a digit beyond the dataset.

\end{document}
